
% This is just an example/guide for you to refer to when submitting manuscripts to Frontiers, it is not mandatory to use Frontiers .cls files nor frontiers.tex  %
% This will only generate the Manuscript, the final article will be typeset by Frontiers after acceptance.   

% When submitting your files, remember to upload this *tex file, the pdf generated with it, the *bib file (if bibliography is not within the *tex) and all the figures.

%%% Version 3.4 Generated 2022/06/14 %%%
%%% You will need to have the following packages installed: datetime, fmtcount, etoolbox, fcprefix, which are normally inlcuded in WinEdt. %%%
%%% In http://www.ctan.org/ you can find the packages and how to install them, if necessary. %%%
%%%  NB logo1.jpg is required in the path in order to correctly compile front page header %%%

\documentclass[utf8]{FrontiersinHarvard} % for articles in journals using the Harvard Referencing Style (Author-Date), for Frontiers Reference Styles by Journal: https://zendesk.frontiersin.org/hc/en-us/articles/360017860337-Frontiers-Reference-Styles-by-Journal

%\documentclass[utf8]{FrontiersinVancouver} % for articles in journals using the Vancouver Reference Style (Numbered), for Frontiers Reference Styles by Journal: https://zendesk.frontiersin.org/hc/en-us/articles/360017860337-Frontiers-Reference-Styles-by-Journal

%\documentclass[utf8]{frontiersinFPHY_FAMS} % Vancouver Reference Style (Numbered) for articles in the journals "Frontiers in Physics" and "Frontiers in Applied Mathematics and Statistics" 

%\setcitestyle{square} % for articles in the journals "Frontiers in Physics" and "Frontiers in Applied Mathematics and Statistics" 

\usepackage{url,hyperref,microtype,subcaption}
%\usepackage{url,hyperref,lineno,microtype,subcaption}
\usepackage[onehalfspacing]{setspace}

%\linenumbers


% Leave a blank line between paragraphs instead of using \\


\def\keyFont{\fontsize{8}{11}\helveticabold }
\def\firstAuthorLast{Vaida {et~al.}} %use et al only if is more than 1 author
\def\Authors{David D Vaida\,$^{1,*}$ and Ryan Jeffrey Farber\,$^{2}$}
% Affiliations should be keyed to the author's name with superscript numbers and be listed as follows: Laboratory, Institute, Department, Organization, City, State abbreviation (USA, Canada, Australia), and Country (without detailed address information such as city zip codes or street names).
% If one of the authors has a change of address, list the new address below the correspondence details using a superscript symbol and use the same symbol to indicate the author in the author list.
\def\Address{$^{1}$ Independent Researcher \\
$^{2}$Department of Physics, Purdue University Fort Wayne, Fort Wayne, IN, USA}
% The Corresponding Author should be marked with an asterisk
% Provide the exact contact address (this time including street name and city zip code) and email of the corresponding author
\def\corrAuthor{David D Vaida}

\def\corrEmail{ddvaida28@gmail.com}




\begin{document}
\onecolumn
\firstpage{1}

\title[Little Red Dots in the JWST Era]{
Little Red Dots: The Assembly of Early Supermassive Black Holes in the JWST Era% Multiwavelength Tests of Compact Accretion at Cosmic Dawn
} 

\author[\firstAuthorLast ]{\Authors} %This field will be automatically populated
\address{} %This field will be automatically populated
\correspondance{} %This field will be automatically populated

\extraAuth{}% If there are more than 1 corresponding author, comment this line and uncomment the next one.
%\extraAuth{corresponding Author2 \\ Laboratory X2, Institute X2, Department X2, Organization X2, Street X2, City X2 , State XX2 (only USA, Canada and Australia), Zip Code2, X2 Country X2, email2@uni2.edu}


\maketitle


\begin{abstract} % currently 208 words, not including keywords
Since the launch of James Webb Space Telescope (JWST) in late 2021, our understanding of high-redshift objects has faced several upheavals. JWST has discovered much more massive galaxies and supermassive black holes (SMBH) than cosmological models had expected. Furthermore, JWST observations have revealed an entirely novel population of high-redshift objects. Characterized by a dominant red rest-frame component and point-like morphology, these ``little red dots'' (LRD) have set off a flurry of observational and theoretical follow-up. The current identity of LRD is highly debated, yet falling into two main scenarios: active galactic nuclei (i.e., SMBH) or compact star-forming regions. If star-forming, LRD would represent the highest stellar densities ever observed. If SMBH, their high Eddington fractions, and already high masses, help elucidate the growth of the most massive SMBH found by JWST in the early Universe ($z \gtrsim4)$. In this mini-review, we present the observational evidence accumulated to date, including sub-millimeter probes of LRD dust masses, constraints on radio and X-ray emission from stacking, and rest-frame ultraviolet \& optical measurements provided by JWST. Furthermore, we highlight how identifying additional LRD that are truly primarily SMBH-driven may help to shed light on the formation of `overly massive' SMBH discovered by JWST within the first billion years since the Big Bang.

\tiny
 \keyFont{ \section{Keywords:} Little Red Dots, Supermassive Black Holes, JWST, AGN, Super-Eddington Accretion, Supermassive Stars, Direct Collapse Black Holes, Primordial Black Holes} %All article types: you may provide up to 8 keywords; at least 5 are mandatory.
\end{abstract}

%% Note: if a line has two percent symbols and a space before the first word, it's one of my comments. 
%% If a line starts with just one percent symbol with no space before the word, I've commented out a line of text 
%% (which will get deleted once its content has been re-worked)
\section{Introduction}
Since the launch of James Webb Space Telescope (JWST; \citealt{gardner2023james,mcelwain2023james}) $\sim$four years ago, JWST has discovered an increasing population of supermassive black holes (SMBH) that challenge the current paradigm of black hole formation and growth (\citealt{jacak2025possible} and see the recent review \citealt{harikane2025early}). 
Specifically, JWST-discovered SMBH are far more massive than expected to be present at such high redshifts ($\sim$billion solar mass black holes are present within the first billion years after the Big Bang), suggesting much more rapid formation than would be possible for Pop III stars accreting at the Eddington limit (\citealt{kiyuna2025super} and see the $\sim$recent reviews \citealt{inayoshi2020assembly,volonteri2021origins,jeon2025physical}). 

Instead, these `overly' massive SMBH may have formed through Eddington-limited accretion from heavy seeds, 
such as
direct collapse black holes from atomic cooling halos and supermassive stars \citep{kiyuna2024sequential,lu2024direct,jeon2025little}, 
or primordial black holes \citep{delos2024structure,riotto2025future,carr2025history}. 
Alternatively, these SMBH may have formed through hierarchical merging of light seed black holes in active galactic nucleus (AGN) disks \citep{vaccaro2024impact}, in dense nuclear or globular star clusters \citep{kritos2025supermassive,lahen2025mergers} or had their growth rates boosted beyond the Bondi rate (and Eddington limit) via fuzzy dark matter soliton cores \citep{chiu2025boosting}. 

Clearly, much uncertainty exists regarding precisely how SMBH grew so rapidly in the early universe. Perhaps the best way to uncover the physical mechanism by which these black holes formed is to find the extension of such a population at high-redshifts that are growing at lower luminosities. Such an intermediate mass population will help to illuminate the seed masses \citep{akins2025cosmos}; moreover, constraints on the accretion rates of such intermediate objects will help understand what seed masses are required for the requisite rapid growth to the observed SMBH masses.

%% Next line: discovery of little red dots: follow the style / use as inspiration the Labbe+ 2025 (2024?) article introduction: phttps://iopscience.iop.org/article/10.3847/1538-4357/ad3551#fnref-apjad3551bib110 
%% Might want to talk about discovery, highlighting the importance of constraining the highly dust obscured portion of the AGN population for constraining growth of early SMBHs (and then talk about the different SMBH formation channels).
Possibly constituting such intermediate mass SMBH in an active state, little red dots (LRD) describe a typically high-redshift population with compact morphology, broad $\sim$1000s km/s H$\alpha$ line widths \citep{kocevski2024rise, zhang2025analysis}, and simultaneously blue ultraviolet (UV) continua yet red optical colors \citep{setton2024little, hviding2025rubies}. Originally, LRD were detected for redshifts $4.2 < z < 5.5$ \citep{matthee2024little}; however, additional JWST observations have uncovered a larger population of LRD extending from redshifts $3 < z < 10$ \citep{labbe2024uncover,graham2025dot}. In addition to a growing population of LRD characterized by JWST, LRD have been investigated by increasingly deep multiwavelength campaigns stretching from radio to X-ray bands.

%%This next line belongs more in the discussion/conclusion
%Several teams now treat LRD as a mixed class where some objects are AGN dominated and others are compact star formers, and they show that mid infrared data are often decisive for separating these possibilities \citep{furtak2025investigating}. 

In this mini-review, we highlight the findings of the most recent LRD observations across the electromagnetic spectrum with a focus on identifying the physical origins of LRD. In section \ref{sec:obs} we contextualize the discovery of LRD with respect to the first few JWST surveys performed. Subsequently, in section \ref{sec:obsConstraining}, we discuss additional multi-wavelength studies that round out the characterization of LRD. 
%we discuss the theoretical work performed to date in interpreting LRD, largely within the AGN framework. 
Synthesizing the observations detailed in prior sections, in section \ref{sec:interp}, we discuss theoretical interpretations of LRD that claim to explain the observed properties of LRD. Finally, we zoom-out and discuss the state of the field of LRD origins more broadly, along with the future directions that can more conclusively end the controversy in section \ref{sec:disc_and_conclude}. Understanding these mechanisms is crucial for accurately interpreting high-redshift observations and for uncovering how the first SMBHs originated.

%% The above is currently 711 words (including the abstract but not really getting the citations word count right; recall the limit is 3000 words for the full article)
\section{LRD Discovery through the Lens of JWST Surveys}
\label{sec:obs}
Some of the first images from JWST, utilizing the Near Infrared Camera \citep[NIRCam;][]{rieke2023performance}, through the Cosmic Evolution Early Release Science program \citep[CEERS;][]{finkelstein2023ceers}, were found to contain very red galaxies (e.g., \citealt{endsley2023jwst,labbe2023population,onoue2023candidate}) implying high-redshift galaxies much more massive than expected from the current cosmological paradigm $\Lambda$ Cold Dark Matter (CDM; \citealt{boylan2023stress,ferrara2023stunning,lovell2023extreme}. 
%CEERS program (Principal Investigator, Finkelstein; Program Identifier, 1345)

Consideration of the longer wavelength properties of these red objects with the Mid-Infrared Instrument \citep[MIRI;][]{argyriou2023jwst} and spectra collected with the Near-Infrared Spectrograph \citep[NIRSpec;][]{bagnasco2007overview} resulted in lower mass estimates, while identifying very broad Balmer line emission and a V-shaped spectral energy distribution (SED): simultaneous blue rest-frame UV and red rest-frame optical colors \citep{carnall2023massive,kocevski2023hidden,barro2024extremely}. 
However, extremely large equivalent widths of the Balmer emission lines suggested a stellar component may contribute a significant fraction to the otherwise clear signature of AGN activity \citep{carnall2023massive}.

%% Next, UNCOVER
Although a small fraction of these red objects were later found to be foreground brown dwarfs, a larger sample of low-luminosity red objects at high-redshift with a characteristic v-shaped SED, point-like morphology, and broad Balmer lines \citep{furtak2023jwst,greene2024uncover,labbe2024uncover} were discovered through the Ultradeep NIRSpec and NIRCam Observations before the Epoch of Reionization \citep[UNCOVER;][]{bezanson2024jwst}) program. 
Combining deep NIRCam photometric and wide-field slitless spectroscopy (WFSS; grism) observations from the EIGER \citep{kashino2023eiger} and FRESCO \citep{oesch2023jwst} surveys, \citet{matthee2024little} coined the term `little red dots' (LRD) to explain the surprisingly abundant population: 
about 1\% of galaxies at redshift $z > 5$ were found to host LRD, while the number density $10^{-5}$ per cMpc$^{-3}$ (co-moving cubic Mega-parsec) is 10-100 times larger than expected by extrapolating the UV luminosity function.


%% {MAYBE DO LATER} Next mention MIRI data / NIRSpec decreases the stellar masses although large dust masses or extremely efficient star formation still seem to be implied.

%% Discuss the RUBIES survey and then our Figure 1 (Hviding 2025 figures 6-8)
Following up on the primarily photometric surveys that incidentally contained little red dots, the Red Unknowns: Bright Infrared Extragalactic Survey (RUBIES; \citealt{de2025rubies}), was designed to specifically investigate LRD through a homogeneous spectroscopic survey. Utilizing NIRSpec and covering cosmic noon ($z \sim$2) and earlier, RUBIES has revealed that LRD show a spectroscopic link between three defining characteristics: broad Balmer lines, a rest-optical point source, and a v-shaped continuum, across 1500 galaxies studied \citep{hviding2025rubies}.

The characteristic bimodal spectral energy distribution was confirmed spectroscopically as one of the defining characteristics of LRD \citep{hviding2025rubies}. In Figure \ref{fig:hviding}, we present three complementary RUBIES diagnostics. First, the $\beta_{\rm UV}$–$\beta_{\rm opt}$ plane shows compact red sources in a tight locus distinct from the parent sample. Second, an Euler diagram quantifies the strong overlap among the three LRD criteria—broad Balmer lines, a dominant rest-optical point source, and a V-shaped continuum. Third, demographics linking redshift to JWST photometric filter flux ratio ratio F356W/F444W and UV absolute magnitude $M_{\rm UV}$ to $L_{\mathrm{H}\alpha}$ indicate that LRD are UV-faint at fixed $L_{\mathrm{H}\alpha}$ yet dominate among the most extreme ${\rm H}\alpha$ emitters. Taken together, these diagnostics tie the characteristic blue-UV/red-optical SED to an unresolved structure and broad Balmer emission, consistent with compact, partly obscured accretion.

%% {MAYBE TO DO} Next JADES, Black thunder??

%% Next the trio of Akins papers...
%% NOTE: Akins+ 2024 is a bit of tentative paper: a single target and marginal detections of He II and Fe X are pretty central to their AGN argument
%% The above is currently 1157 words (including the abstract but not really getting the citations word count right; recall the limit is 3000 words for the full article)
\section{Observational Constraints on LRD}
\label{sec:obsConstraining}
%\subsection{Rest-frame UV}
Studying the rest-frame UV of LRD, detections of the high-ionization coronal line [Fe\small{X}] strongly suggest an AGN nature \citep{kocevski2023hidden,furtak2024high}. 
On the other hand, non-detections of He\small{II} emission lines in the far UV (FUV) and Mg\small{II} are more consistent with photoionization by massive stars \citep{akins2024strong}. 
Additionally, a Balmer break, observed for many LRD, typically suggests an aged stellar population, though it is not universal \citep{akins2024strong}. 
On the other hand, detections of strong C \textsc{iv}, marginal He \textsc{ii} and [Fe \textsc{x}], together with broad H$\alpha$ all combined strongly support an AGN interpretation \citep{inayoshi2025extremely}. 
%(However, the CIV is consistent with local dwarfs with high SFR and low metallicity). 

%% NOTE: Akins+ 2025 tentative...is indeed a tentative paper - 3 targets for one thing.
%\citet{akins2025tentative} utilize ALMA to measure the amount of CO and [C$\small{I}$] for 3 LRD at redshifts, 3,4, and 7. They place upper limits on CO and find only a tentative detection of [C$\small{I}$] in one of the brightest LRD, suggesting LRD are gas-poor compared to high-redshift quasars. Furthermore, the narrow full-width at half-maximum (FWHM) of the [C$\small{I}$] line suggests a dynamical mass an order of magnitude smaller than the stellar mass found from SED fitting, suggesting the observed Balmer breaks are non-stellar in origin. The dynamical mass limit is also lower than the local $M_{\rm BH}/M_{\rm dyn}$ relation.

%\subsection{Far-IR \& Sub-millimeter: Dust Constraints}
As early as 2024, it was reasoned that (most) LRD dust masses must be quite limited $< 10^{4-5}\,M_\odot$ due to their compact sizes yet relatively low mass available for attenuation \citep{casey2024dust}. Indeed, deep NOEMA observations of $z > 7$ LRD \citep{xiao2025no} as well as multi-band Atacama Large Millimeter/Submillimeter Array \citep[ALMA;][]{wootten2009atacama} observations of a couple of the brightest LRD known resulted in stringent non-detections \citep{setton2025confirmed}. Observations of a larger sample ($\sim 60$) LRD with ALMA in the 1.3 mm continuum band similarly resulted in no detections, even with stacking, suggesting LRD either have modest dust reservoirs  ($A_v \sim 2-4$) or otherwise very dense gas $10^9\,\rm{cm}^{-3}$ causes the obscuration. The paucity of dust in LRD extends to a broader range of wavelengths in the far-IR and sub-mm: LRD observations in the rest-frame IR in Spitzer/MIPS 24 $\mu$m, JCMT/SCUBA-2 850 $\mu$m, ALMA 1.2 \& 2.0 mm, and MeerKAT/VLA 1.3 \& 3.0 GHz all resulted in non-detections \citep{bao2025investigating}.

Consistent with non-detections at redder wavelengths, \citet{brooks2025here} investigated the dust attenuation of Balmer narrow-line emission and broad-line signatures in 29 $z > 3.5$ AGN selected from JWST surveys (CEERS, JADES, and RUBIES). The narrow-line emission is consistent with no dust attenuation, which \citet{brooks2025here} consider to be consistent with star formation at larger scales in the galaxy. They place a lower limit on the dust attenuation $A_v > 3.6$ due to non-detection ($<3\sigma$) of the broad H$\beta$ line even after stacking 25 of their sources, which they say is consistent with a dust-obscured AGN. 


%They suggest the red-colors of LRD are due to dust attenuation from AGN emission while the blue UV colors (traced by the narrow line emission) is predominantly related to star formation in the host galaxy (narrow-line regions have scales of 3-10 kpc). Since dust obscuration by a canonical AGN torus is inconsistent with non-detections in mid-IR (though those non-detections were for targets in the GOODS-S survey which doesn't overlap well with Brooks targets), they propose the dust is still in the galactic nucleus (compact compared to star-forming scales) but farther from the AGN than a canonical torus or consistent with polar dust which would be colder than the torus. The presence of strong Balmer absorption lines in two RUBIES sources suggests very high neutral hydrogen densities. They say the difference in dust attenuation between narrow and broad Balmer line components favors an AGN model; a star-formation model would expect to have more continuous velocity and attenuation distributions.

%% Bao comparison to super-Eddington lower redshift objects
%in compared blue hot dust-obscured galaxies to LRD. They find that the v-shaped SED occurs at redder wavelengths for blue hot dust-obscured galaxies, suggesting that LRD have less dust extinction or at least, no dust heated by AGN activity, 
%They argue against the blue excess in LRD being due to AGN scattered light too. Overall, LRD are different than blue hot dust-obscured galaxies, which are thought to be powered by super-Eddington accretion of the SMBH after major mergers.

%% This one can definitely stay - it should go ~last.
Multi-wavelength observations resulting in non-detections have extended to an increasingly larger sample of LRD. Studying 434 LRD in the JWST COSMOS-Web survey, \citet{akins2025cosmos} find 
%that their brightest objects require AGN contributions to be consistent with the maximum surface densities observed in the local universe \citep{hopkins2010maximum}. On the other hand, they found 
no detections, even when stacking, in X-ray, mid-IR, far-IR/submillimeter, and radio bands. 
Their non-detections are consistent with prior work of smaller samples that found LRD are not detected, or quite weakly emitting, in X-ray \citep{ananna2024x,yue2024stacking,maiolino2025jwst} and in radio \citep{mazzolari2024radio,perger2025deep,gloudemans2025another}.


%\citet{carranzaEscudero2025lonely} analyzed 124 LRD from the CEERS, NEP-TDF, JADES, and JEMS surveys. North Ecliptic Pole Time Domain Field (NEP-TDF; Windhorst et al. 2022) survey. They found that while the SED fit improved with inclusion of an AGN component, their Bayesian information criterion suggested the improved $\chi^2$ is due to overfitting; therefore, models without an AGN are more appropriate for characterizing LRD, especially when MIRI data is available. They find that LRD are in less dense environments compared to galaxies at all masses and similar redshifts (two populations according to the Kolmogorov-Smirnov test with p-values < 0.05 (circa 0.01)). They suggest most LRD are compact galaxies or with active star formation. Using abundance-matching, they find no evolution in the stellar-mass halo-mass relationship from redshifts 3-10.


%Counter to \citet{carranzaEscudero2025lonely} finding LRD are in less dense environments, is the assertion LRD are found in over-densities [intro of Furtak 2025 investigating paper]  (e.g. Schindler et al. 2024; Matthee et al. 2024b)

%While the lack of X-rays \& radio would otherwise be strong evidence against an AGN-nature of LRD, given the high-redshift nature of LRD, there could be high gas column densities causing Compton-thick absorption of the X-rays and the sources to be radio-quiet. While the far-IR constraints rule out hot dust in a prototypical AGN torus near the central SMBH, there could be substantial warm dust below their detection limits \citep[see also ][]{akins2025tentative}.


%% We should probably start by mentioning the main surveys that have been performed (black thunder, rubies, jades, etc.) and summarize their results on the properties of LRD. Then, I think it'd be useful to discuss the lack of X-rays and variability, suggesting a stellar origin. Last, we discuss the observations by ALMA etc of dust suggesting X-rays are obscured, leaving the question fairly open.

%Follow up catalogs expanded the set with lensing assisted discoveries such as the “Ruby” and with systematic RUBIES style selections that emphasized size, color, and line criteria \citep{de2025remarkable}. 

%% I'll need to massage this
%Originally, most LRD were discovered photometrically using JWST's NIRCAM \citep{barro2024extremely}.
%, which show a variety of features, from high equivalent widths (i.e., extinction due to dust) explaining the red colors for some sources \citep{larson2023ceers}, whereas other LRD exhibit intrinsically red colors though occasionally with spectral breaks.

%Follow-up of high-redshift $5 < z < 9$ morphologically point-like, extremely red objects from the Cosmic Evolution Early Release Science Survey \citep[CEERS;][]{finkelstein2023ceers}, confirmed the defining features of LRD. 
%\citet{barro2024extremely} found rest-frame blue ultraviolet (UV) colors in these objects, with JWST's Near Infrared Camera \citep[NIRCam;][]{rieke2023performance}), while observations with the Mid-Infrared Instrument \citep[MIRI;][]{argyriou2023jwst} indicated a steep slope. 

%% Discuss the Furtak figure probably in the interpretation section showing clear AGN evidence due to the broad line variability.




%% lacking mid-IR (no hot dust)
%(e.g. Williams et al. 2024; Akins et al. 2024;





%% from Casey intro, low-z LRD
%low-z (z < 3) LRD –both compact and red – seem to be far more rare (Lin et al. 2025a; Bisigello et al. 2025; Ma et al. 2025).

%% interpretation -- cannot be stars bc Balmer break strength is too high (Casey intro)
%This idea is supported by recent observations of LRD with extreme Balmer break strengths, exceeding the maximal expected from a stellar population (Naidu et al. 2025; Rusakov et al. 2025; de Graaff et al. 2025; Ji et al. 2025; Taylor et al. 2025).




% %% AGN evidence
% \citet{d2025blackthunder} follow-up on a source they previously found time-variable equivalent width in broad emission lines (confirmed its AGN nature) from the BlackTHUNDER survey of NIRSpec/IFS data. They find a narrow-line H$\alpha$ component, confirming a connection between the strong Balmer break and absorption line features. With a grating spectrometer they find a black hol to dynamical mass ratio of 2-40\%: either super-Eddington accretion or $\sim$no host galaxy for the min/max cases. From the Balmer decrement in the narrow lines, they find low dust extinction $A_v \sim1.4$, consistent with low dust mass \citep{casey2025upper}.

%% copied from c-e conclusion
%The SHMR is greater than the maximum SHMR at z = 0 for four redshift bins. Because of this, we expect that either all massive halos host LRD, in which case we would expect LRD to be found in pairs, or AGN provide a large fraction of the light from LRD.

%% important articles arguing LRD are stars not AGN (sourced from carranza-escudero 2025)
%Perez-Gonzalez et al. (2024) find that the red optical and near infrared (NIR) data fit an obscured accretion disk, but stellar-dominated (often starburst) models provide an even better fit. Williams et al. (2023) used MIRI+ALMA vs ALMA only

% %% about LRD having overmassive BH compared to host galaxy stellar mass (BlackTHUNDER strikes twice intro; D'Eugenio 2025
% over-massive relative to their hosts’ stellar mass (e.g. Kocevski et al. 2023; Übler et al. 2023; Harikane et al. 2023; Kokorev et al. 2023; Maiolino et al. 2024b).
% Also \citep{d2025blackthunder}. Also \citep{durodola2025exploring} I guess.
% While $M_*$ measurements in AGN-dominated systems remain uncertain, observations of faint SMBHs with low accretion rates support this trend (Carnall et al. 2023; Juodžbalis et al. 2024b).
% %% Related but from Durodola+ 2025
% Interestingly, despite the deviation from the local $M_{\rm bh}-M_*$ relation, the local relations $M_{bh}-\sigma$ and $M_{\rm bh}-M_{\rm dyn}$ appear to hold (R. Maiolino et al. 2024)
% %% More overmassive (10-100x) sources from Durodola+ 2025
% F. Pacucci et al. 2023.
% %% The whole Graham paper is about the morphology-dependence of the M_bh - M_star relation and say that LRD from 2023-2024 were below the relation for nuclear star clusters but similar to ultracompact dwarfs and primeval lenticulars, whereas 2025 LRD may be similar to the relation for nuclear star clusters, green peas, and UCD galaxies.

% %% argument against strong Balmer breaks being due to aged stars in LRD  (BlackTHUNDER strikes twice intro; D'Eugenio 2025
%  strong Balmer breaks; if these spectral features arise from evolved stellar populations (a few hundred Myr and older), they would imply the existence of massive, old galaxies 600 Myr after the Big Bang (Wang et al. 2024a,b) with extremely high surface densities (Baggen et al. 2024; Ma et al. 2024). Alternatively, these Balmer breaks could arise from dense gas absorption (Inayoshi \& Maiolino 2025) along the line of sight to the SMBH.

% %% argument for AGN over stars interpretation (BlackTHUNDER strikes twice intro; D'Eugenio 2025
% The discovery of Abell2744-QSO1 (Furtak et al. 2023, 2024) enabled us to break this degeneracy. The continuum of this source cannot be well reproduced by stellar-population spectra alone (Ma et al. 2024, at least not without invoking ad hoc attenuation laws). Subsequent JWST observations have revealed that the equivalent widths (EW) of the broad lines are time variable (Ji et al. 2025, hereafter: J25; Furtak et al. 2025), ruling out a stellar origin (e.g., Baggen et al. 2024; Kokubo \& Harikane 2024). Furthermore, time variability of the rest-frame optical continuum (J25) implies an AGN origin (but see Furtak et al. 2025 for a different view)
% Stars produce Gaussian broad lines whereas \citet{d2025blackthunder} find non-Gaussian (requires two Gaussian components).

% %% LRD with broad Balmer absoprtion (BlackTHUNDER strikes twice intro; D'Eugenio 2025
% Mattee et al. 2024; Wang et al. 2024a,b; Juodžbalis et al. 2024a; Taylor et al. 2024; Labbe et al. 2025).

% %% no variability of LRD (c-e 2025 paper intro)
% For example, Kokubo \& Harikane (2024) report that LRD do not exhibit the typical photometric variability associated with standard AGNs. However, Kokubo \& Harikane (2024) suggest that this could be the result of intrinsically non-variable AGN or the dominance of AGN emission through scattering. Zhang et al. 2025  also find that LRD do not usually display strong
% photometric variability.
% %% and from Furtak 2025 investigating
% (Kokubo \& Harikane 2024; Zhang et al. 2024; Tee et al. 2025; Ji et al. 2025). The investigated samples show mostly no variability in several epochs of broad-band photometry, both in the rest-frame ultraviolet (UV) and optical (Kokubo \& Harikane 2024; Zhang et al. 2024; Tee et al. 2025) with only a small subsample of LRD displaying potential variability (Zhang et al. 2024). However, since JWST has only been operational for a little over two years, time-domain data are still sparse. 

% %% copied from brooks, need to re-work: (but important about X-rays for LRD with AGN interpretation)
% Analyses of two X-ray-confirmed AGN, at z = 3.1 and z = 4.66 in B. Wang et al. (2025) and D. D. Kocevski et al. (2024), find that their X-ray emission is consistent with a dust-reddened AGN and that they could be lower-redshift analogs of the LRD population.

% \citet{bellovary2025little} propose a tidal-disruption event explanation of LRD. Their TDE model explains the broad H$\alpha$ emission, which is then independent of the AGN mass (therefore, the SMBH are not necessarily over-massive relative to their galactic hosts), H$\alpha$ absorption (due to tidal streams of disrupted stars), weak X-ray emission due to super-Eddington accretion combined with viewing angle effects (though the viewing angle explanation seems inconsistent with the number density of LRD and the blue hotdog paper suggests LRD are not accreting super-Eddington). They model the red-component of the SED to come from stars, and they point out a top-heavy initial mass function (IMF) can alleviate concerns of the stellar mass being higher than supported by $\Lambda$CDM. Furthermore, the IMF has been shown to vary with density and metallicity (Marks et al. 2012). Ultimately, \citet{bellovary2025little} suggests some fraction of LRD may be explained as TDE since some LRD show clear AGN signatures. However, the TDE hyptohesis can be tested by finding variability in the rest-frame UV while the optical-NIR remains fairly constant, which may be observed by JWST within the next five-ten years.

% %% Future work (pasted from Akins 2025 cosmos)
% Future large spectroscopic surveys will be needed to conduct a comprehensive census of LRD and to elucidate their true nature at a population level, and deeper observations in the X-ray, mid/far-IR, and radio will be needed to determine their true properties and draw connections to the full AGN population.

% %% Interpretation
% \citet{cenci2025little} present preliminary runs from the MELI$\odot$RA cosmological hydro simulations that include accurate treatment of the formation of direct-collapse black holes by handling Lyman-Werner radiation and mass inflow in target halos. They find newly formed DCBHs drop off for $z < 6$ due to reduced inflows. When the DCBH is formed, there is a gas compaction event, followed by $\sim$200\,Myr of intense luminosity, and afterwards the formation of the first Pop III stars. 
% Their consideration of very low metallicities for DCBH formation $10^{-4-5}\,Z_\odot$ appears to be consistent with the low mass of dust found in LRD \citep{casey2025upper}.
% I'm confused about Pop III stars formation $\sim$after the LRD in their model though given the attribution of narrow-line Balmer emission to stars \citep{brooks2025here}.

% %% Evidence for heavy seeds
% \citet{durodola2025exploring} consider 19 LRD from JADES with NIRCam and MIRI photometry. They find one with a 2\% AGN contribution to the luminosity and one completely star-forming. The rest have 20-70\% (or 40-85\% according to the conclusion section) contributions from AGN which are significantly overmassive compared to their hosts (1-4 dex above the local $M_{\rm bh}-M_{*}$ relation.

% %% cenci abstract
% LRD clearly emerge as a distinct population at redshifts 5-8 and clearly decline for redshifts less than five.

% %% cenci intro (but also casey intro)
% LRD constitute a few percent of the total galaxies at high-redshifts (non-negliglible; Pizzati+ 2025 and also Arita+ 2025, Matthee+ 2025) and are $\sim$low mass halos $\sim10^{11} M_\odot$.
% %% furtak+ 2025 investigating
% Though 20-30\% of the AGN population is LRD (so pretty high)  (e.g. Kocevski et al. 2024)

% %% cenci intro (weak x-ray AGN interp -- high neutral gas column densities at super-Edd accretion)
% (e.g. Inayoshi, Kimura \& Noda 2024; Lambrides et al. 2024; Pacucci \& Narayan 2024; Madau 2025)
% %% cenci intro; theoretical interp as heavy seed BH growth
% see Inayoshi 2025; Inayoshi \& Maiolino 2025; or supermassive star progenitors: Begelman \& Dexter 2025

% %% Interpretation
% \citet{berton2025similar} review the properties of narrow-line seyfert type-1 (NLS1) AGN to LRD. NLS1 are low-redshift AGN with broad H$\alpha$ emission lines $\sim$2000 km/s, low SMBH masses $\sim10^{5-7}M_{\odot}$ which sound reminiscent of LRD. However, \citet{berton2025similar} points out H$\alpha$ absorption features observed in LRD are more reminiscent of P-Cygni profiles, and are not observed in NLS1. Regarding LRD being X-ray weak, \citet{berton2025similar} point out that NLS1 have a steep X-ray slope and Chandra measures a harder X-ray component in high-redshift sources, consistent with weak or non-detections. NLS1 also have weaker X-ray emission compared to their UV and the outflow properties of high-z AGN (e.g., LRD) may further suppress X-ray emission.

% %% Interpretation (support for an AGN picture) - comes from Akins 2025 cosmos
% Reconciling these high masses with the pointlike morphology requires rapid starbursts in compact cores (see, e.g., J. F. W. Baggen et al. 2023), and implies a high stellar mass density in the early Universe, which, for some time, appeared to be in tension with expectations for the halo mass function in $\Lambda$ cold dark matter (CDM; M. BoylanKolchin 2023). While more recent analysis has, in some cases, brought down the exceedingly large stellar masses, relieving the possible tension with $\Lambda$CDM (e.g., R. Endsley et al. 2023b; K. Chworowsky et al. 2024; D. D. Kocevski et al. 2023), in others, the high redshifts and high stellar masses have been reaffirmed by spectroscopic observations (e.g., B. Wang et al. 2024a). 

% %% Interpretation
% Weak observed H$\beta$ (e.g., \citealt{akins2024strong}) could be due to super-Eddington accretion.  Pacucci \& Narayan 2024; Lambrides et al. 2024) [see Akins 2024 paper for the citations] , or just dust attenuation. However, the high-ionization lines marginally detected are inconsistent with high-Eddington accretion (radiatively inefficient). 


% %% Interpretation (not an obs paper)
% \citet{begelman2025little} argue that LRD are quasi-stars: hypothetical objects that occur after a supermassive star (SMS) collapses to a seed black hole, but during a phase in which the massive envelope of the remnant stars is still convective: the seed SMBH accretes at a super-Eddington rate until the envelope is depleted. According to \citet{begelman2025little}, in this case all galaxies undergo a brief LRD period of $\sim$10\,Myr. There model can be tested by observing the evolutionary sequence from early-type LRD, when the accretion rate was even higher, to the presently observed extended late-type LRD stage. Or less ambiguously given the diversity of LRD, observations of SMS precursors to LRD would help with the identification. They explain that quasi-stars have no X-ray emission due to high electron column densities.

% %% Related black hole star papers (sourced from carranza-escudero 2025 intro)
% Theories discussing ”black hole stars” (Naidu et al. 2025; de Graaff et al. 2025) or hyperdense gas cocoons surrounding smaller black holes also exist (Inayoshi \& Maiolino 2025), with early ALMA results providing upper mass limits that are in line with this (Casey et al. 2025).


% %% Interpretation
% Whereas different stellar populations in different regions of a galaxy can produce both red and blue UV-optical colors,
% \citet{barro2024extremely} point out that the unresolved point-like source together with the bimodal UV-optical color distribution is difficult to explain with a stellar model, favoring an AGN scenario. 
% %% Below needs work.
% On the other hand, they cannot distinguish between AGN-alone or QSO + blue star SED models (degenerate), requiring redder wavelength coverage (e.g., still MIRI).

% %% This is related to interpretation pretty closely and is arguably the most useful in characterizing the nature of LRD as AGN vs star-forming
% Usually, observations in the rest-frame far UV (FUV) and optical (covering the H$\alpha$ and H$\beta$ lines) would be sufficient to characterize the star formation rates\citep[e.g.,][]{bell2001comparison}, clarifying whether the cause of the bright emission in LRD is due to stars or quasi-stellar objects (i.e., AGN) through the BPT diagram (which compares flux ratios of [N \small{II}]/H$\alpha$ to [O \small{III}]/H$\beta$; \citealt{baldwin1981classification,kewley2001theoretical,kewley2013cosmic}). 
% However, regions with high dust extinction under-predict star formation rates based on FUV and Hydrogen Balmer lines alone \citep{calzetti2001dust,buat2002star}. 
% To characterize the dust mass (hence the magnitude of extinction), and independently determine the star formation rate, rest-frame far-infrared (IR) measurements are required \citep{garn2010obscured}.
% Therefore, coverage in the rest-frame far-infrared with the Atacama Large Millimeter/Submillimeter Array \citep[ALMA;][]{wootten2009atacama} is necessary for robust characterization of the star formation rate, even in regions with high amounts of extinction by dust in the rest-frame far UV \citep[FUV;][]{khusanova2021alpine,algera2023alma}. 
% ALMA observations in combination with Herschel \citep{pilbratt2010herschel} of $z \sim 2.5$ LRD suggest very dusty environments with most objects undergoing intense episodes of star formation (though one appears to be in a quiescent post-starburst state \citealt{cooper2025rubies}).

% However, further characterization in the mid-infrared with MIRI \citep{barro2024extremely}, ...
% Further multiwavelength comparison utilizing X-ray ...

% from Cooper
%The NIRSpec microshutter array (MSA) survey, Red Unknowns: Bright Infrared Extragalactic Survey (RUBIES; GO#4233, PIs: A. de Graaff and G. Brammer; A. de Graaff et al. 2024) is designed to achieve spectroscopic completeness for bright, red objects at cosmic noon and earlier. Target selection and mask design is optimized to return high completeness in particular for the most extreme sources in this color space, which includes DSFGs. At the redshift regime where most DSFGs sit (1 < z < 3), RUBIES covers multiple Paschen and Balmer lines, in addition to other species emerging from nebular regions. Joined with far-infrared (FIR) constraints from the Atacama Large Millimeter/submillimeter Array (ALMA), these data enable robust measurements of nebular, stellar, and dust properties, as well as direct constraints on ongoing star formation activity and history.


% % Evidence for AGN
% Observational support for accretion comes first from direct line diagnostics. Spectroscopy has revealed broad $H\alpha$ and in some cases $H\beta$ line widths consistent with a compact broad line region, and repeat observations in a lensed system show equivalent width variability on $\sim$year timescales consistent with AGN \citep{d2025blackthunder,ji2025blackthunder}. That is,  at least a fraction of LRD host accreting black holes rather than pure stellar continua. 
% % Evidence for stellar origin
% Photometry that combines NIRCam with MIRI often prefers composite stellar plus AGN fits for objects with red rest-optical colors and blue UV slopes, and the addition of mid infrared points reduces the fraction of sources explainable by AGN only templates, which argues for a population mix rather than a single class \citep{labbe2024uncover}. 
% Host properties add a final observational thread. Some LRD are extremely compact with low dynamical or stellar masses, and when single epoch line methods are used the implied black hole to host ratios can exceed local relations, a pattern that is easier to reconcile if early accretion sometimes outpaces stellar assembly \citep{naidu2025black}.
% Lensing aided measurements strengthen this point for at least one Ruby-like system where the dynamical mass is low enough to strain purely stellar Balmer break interpretations \citep{de2025remarkable}. Altogether these observations indicate that a meaningful subset of LRD is consistent with actively growing black holes even when classical X ray signatures are weak.
% % X-ray and radio suggest stellar or high obscuration
% Stacking results provide a second line of evidence. X ray stacks and radio stacks are predominantly non detections yet remain consistent with low luminosity or heavily obscured AGN that would have large bolometric corrections, and one secure AGN in the broader candidate set demonstrates that the selection can recover bona fide accretors when conditions are favorable \citep{hainline2025investigation}. 
% Population level context supports these inferences. The faint end X ray luminosity function at z greater than four suggests a reservoir of accretors with weak coronae, which lines up with the X ray quiet nature of many LRD \citep{barlow2025measurements}. 
% X-ray diagnostics reviews provide the physical basis for how thick inner flows and winds mute the hard emission while leaving strong rest-optical output, which matches several LRD spectral energy distributions \citep[SED;]{kara2025supermassive}. 

%% The above is currently 1552 words (including the abstract but not really getting the citations word count right; recall the limit is 3000 words for the full article)
\section{Theoretical Interpretations}
\label{sec:interp}
We next consider how LRD spectral properties arise and why AGN and stellar interpretations can both appear.
Forward modeling shows that an AGN continuum filtered by dense nearly dust free gas can generate a strong Balmer break and broad Balmer lines while leaving the UV relatively blue, a combination that matches several spectroscopic LRD with tentative variability present \citep[2.6$\sigma$;][]{d2025blackthunder}. 
In particular, the inclusion of mid-IR and redder data (from MIRI \& ALMA) has led to mixed results. 
While the red optical and near-IR data can be fit by an obscured accretion disk with scattered AGN light \citep{lambrides2024case,pacucci2024mildly,madau2025chasing}, stellar-dominated (even starburst) models provide improved fits \citep{perez2024nature,carranzaEscudero2025lonely}. 
Moreover, the concerns that a stellar nature of LRD is inconsistent with the maximal brightness surface density in the local universe \citep{hopkins2010maximum}, is alleviated due to MIRI and ALMA observations decreasing the requisite stellar mass \citep{williams2024galaxies}. 
On the other hand, recent analysis of LRD from the
JWST Advanced Deep Extragalactic Survey \citep[JADES;][]{eisenstein2023overview}
with NIRCam and MIRI photometry suggest one LRD with a 2\% AGN contribution to the luminosity, one completely star-forming, and the rest with 20-70\% contributions from AGN \citep{durodola2025exploring}. 
Similarly, recent analysis of MIRI data increases the fit for stellar components, reducing the number of AGN-only solutions without eliminating AGN altogether \citep{furtak2025investigating}. 

To explain the paucity of radio, X-ray, and far-IR/sub-mm emission, modeling work has extended beyond the standard models of AGN. 
Quasi star or black hole envelope models predict compact very red sources with strong Balmer features during specific evolutionary phases that overlap LRD selections, suggesting some LRD may represent heavy seeds in formation \citep{begelman2025little,durodola2025exploring}. 
Black hole envelopes embedded in dense gas can also reproduce red continua and strong lines in a way that mimics faint AGN \citep{kido2025black,naidu2025black}. 
On the other hand, simulations of collapsing clusters predict enhanced tidal disruption activity and luminous transients that can generate broad lines and variable continua in compact systems, offering another channel that suggests a non-standard AGN origin of LRD with rapid black hole growth in cluster cores \citep{bellovary2025little}. 
Dusty inflow models explain very red rest-optical colors and broad absorption and predict that mid-IR observations should be constraining \citep{li2025little}. 
%Tools built for broadband quasar SEDs allow stress tests that compare AGN dominated and compact starburst templates across the NIRCam to MIRI baseline, which gives a practical route to refine classifications in future samples \citep{saccheo2025hyperion}. 
Simulations of high redshift environments indicate that compact AGN can arise in gas rich protoclusters where repeated fueling and multiphase outflows sculpt line profiles and continuum shape
%, which provides a physical context for LRD that show strong winds and complex absorption 
\citep{kannan2025introducing}. 
Observations of multiphase outflows around $z \sim 5$ quasars show that such winds can be common and energetic, which supports the idea that some LRD reflect wind-bearing accretion states rather than aged stellar populations \citep{brazzini2025multi}. These models collectively predict that modest changes in geometry, column density, and recent fueling history can toggle an object between apparently AGN-like and stellar-like diagnostics, which explains why LRD form a mixed class.

%% strong AGN evidence
%% DAVID TO DO - tune this paragraph up and discuss the Furtak figure.
Considering a multiply imaged LRD at redshift 7 from UNCOVER, \citet{furtak2025investigating} investigated time variability, leveraging lensing time delays of 22 years (rest-frame $\sim3\,$yr). Finding significant variability in H$\alpha$ and H$\beta$ lines, \citet{furtak2023jwst} confirmed the AGN nature of this LRD. We show \citet{furtak2025investigating} their figure 1 in Figure \ref{fig:furtak}. Clearly, the H$\alpha$ and H$\beta$ line fluxes have evolved significantly from epoch 7 (marked in blue) to epoch 8 (marked in green). This constitutes strong AGN evidence for several reasons. First, the Balmer lines are very broad (FWHM ~ a few $10^3kms^{-1}$), which requires the high velocities of gas near a black hole and cannot be produced by normal star-forming H \textsc{ii} regions. The line fluxes and profiles vary on rest-frame $\sim$year timescales, the expected behavior of a compact broad-line region responding to a changing ionizing source. Since the spectra are normalized and the multiple images are separated by known lensing delays, the differences cannot be due to calibration or lensing and must be intrinsic. Photometric non-variability does not contradict this result, since AGN continuum variability is stochastic and can be modest over the sampled interval. 

%% The above is currently 2160 words (including the abstract but not really getting the citations word count right; recall the limit is 3000 words for the full article)
\section{Discussion and Conclusion}
\label{sec:disc_and_conclude}
%Limitations and open questions follow from selection, diagnostics, and environment. 
Despite recent progress, many caveats surround the work performed to-date attempting to identify the nature of individual LRD. Many LRD are X-ray and radio quiet even in stacks, which could be consistent with heavy obscuration and low coronal efficiency (as expected from super-Eddington accretion models; \citealt{secunda2025little}) or purely stellar solutions; stellar templates fit a large fraction or LRD with mid-IR data included \citep{hainline2025investigation}. 
Radio and X-ray constraints specific to compact high redshift sources show that non detections do not rule out accretion for the luminosities expected from thick trapping dominated flows, which raises the need for sensitive follow up tuned to low efficiency spectra rather than relying on shallow stacks \citep{latif2025radio}. Deeper observations in radio bands may uncover a tail of detections if even a minority of LRD host jets or compact cores, which provides a near-term observational test of the mixed population picture \citep{latif2025synergies}. 

%Catalog and SED work confirms that the AGN fraction is sensitive to mid infrared depth and to the exact template set, so population fractions remain uncertain until deeper MIRI and longer wavelength data are common across fields \citep{furtak2025investigating}. 
Completeness studies for obscured narrow line AGN show that color cuts and emission line thresholds can miss faint or partially obscured accretors, which implies that some LRD-like systems may be absent from current AGN catalogs and that selection bias can skew interpretations \citep{bouwens2025beneath,scholtz2025jades}. 
%UV selected studies of massive quasar hosts argue that some AGN at high redshift are UV faint relative to their mass, which complicates color based selections and suggests that LRD compilations can be incomplete in ways that correlate with SED shape \citep{bouwens2025beneath}. 

Ultimately, LRD likely comprise both genuine rapidly growing black holes and compact star-forming systems, and resolving their roles in early black hole demographics will require deeper mid-IR coverage, higher signal to noise spectroscopy that targets both Balmer and high ionization UV lines, and radio and X-ray follow up that is matched to the efficiencies expected for thick discs rather than thin ones \citep{akins2025tentative}.

In particular, measuring time variability would provide `smoking-gun' evidence for an AGN nature of individual LRD. As most LRD are not multiply imaged as in the target of \citet{furtak2025investigating}, collecting additional epochs of LRD photometric and spectroscopic measurements extending several years will be necessary. The upcoming \texttt{TWINKLE} campaign (JWST Proposal Cycle 4, ID 7404PI, PI: Naidu) will be the next possibility to find variability in the Balmer emission line fluxes.
While expensive, deeper observations of truly AGN-generating LRD will be crucial to better understand the assembly of supermassive black holes within the first billion years from the Big Bang.

%Thanks to the gravitational time delays, this source is ideally suited for further study of its variability and could make for a prime target for a large-scale reverberation mapping campaign in the future. With that, one could potentially measure the black-hole mass of an AGN at z = 7 directly and, for the first time, calibrate the scaling relations used to measure black-hole masses beyond the near-by Universe (e.g. Liu et al. 2024). Thanks to SL time delays, multiply-imaged, clusterlensed AGN like A2744-QSO1 may open a door to achieve this on a reasonable timescale out to high-redshifts and high blackhole masses (e.g. Golubchik et al. 2024)

%The environment of LRD remains unsettled. 
%Samples from RUBIES suggest LRD are preferentially found in low-density environments  while protocluster simulations predict that compact AGN should be common in overdensities, and reconciling these views will require larger spectroscopic samples with uniform selection across environments. 


\section*{Conflict of Interest Statement}
%All financial, commercial or other relationships that might be perceived by the academic community as representing a potential conflict of interest must be disclosed. If no such relationship exists, authors will be asked to confirm the following statement: 

The authors declare that the research was conducted in the absence of any commercial or financial relationships that could be construed as a potential conflict of interest.

\section*{Author Contributions}
DDV: Writing -- original draft, review and editing. RJF: Writing -- review and editing.

\section*{Acknowledgments}
We thank Lukas Furtak and Adi Zitrin for their permission to re-use their figure as shown in our Figure \ref{fig:furtak}.
We thank Raphael Hviding for their permission to re-use their figure as shown in our Figure \ref{fig:hviding}.
DDV thanks his family for their belief and support throughout the writing process.

\bibliographystyle{Frontiers-Harvard} %  Many Frontiers journals use the Harvard referencing system (Author-date), to find the style and resources for the journal you are submitting to: https://zendesk.frontiersin.org/hc/en-us/articles/360017860337-Frontiers-Reference-Styles-by-Journal. For Humanities and Social Sciences articles please include page numbers in the in-text citations 
%\bibliographystyle{Frontiers-Vancouver} % Many Frontiers journals use the numbered referencing system, to find the style and resources for the journal you are submitting to: https://zendesk.frontiersin.org/hc/en-us/articles/360017860337-Frontiers-Reference-Styles-by-Journal
\bibliography{test}

%%% Make sure to upload the bib file along with the tex file and PDF
%%% Please see the test.bib file for some examples of references

\section*{Figure captions}

%%% Please be aware that for original research articles we only permit a combined number of 15 figures and tables, one figure with multiple subfigures will count as only one figure.
%%% Use this if adding the figures directly in the mansucript, if so, please remember to also upload the files when submitting your article
%%% There is no need for adding the file termination, as long as you indicate where the file is saved. In the examples below the files (logo1.eps and logos.eps) are in the Frontiers LaTeX folder
%%% If using *.tif files convert them to .jpg or .png
%%%  NB logo1.eps is required in the path in order to correctly compile front page header %%%

\begin{figure}[ht]
\begin{center}
\includegraphics[width=\textwidth]{HvidingComposite.png}% This is a *.eps file
\end{center}
\caption{
Adapted from \citet{hviding2025rubies}, their Figures 6, 7, and 8 to composite the RUBIES diagnostics for Little Red Dots. 
In the upper left `corner' plot and upper-middle plot, the $\beta_{\text{UV}}$--$\beta_{\text{opt}}$ maps show where sources with broad Balmer lines, unresolved rest-optical point sources, and V-shaped continua sit relative to all $z_{\text{spec}} > 3.1$ objects. 
In the upper right, the two heatmaps respectively show the photometric flux in two filters for high vs. low redshift sources (threshold $z = 5$) as a function of spectroscopic redshift, and the spectroscopically derived UV absolute magnitude versus the H$\alpha$ luminosity. Those two heatmaps show LRD are UV-faint at fixed $L_{\text{H}\alpha}$ and dominate the most H$\alpha$-luminous objects at fixed $M_{\text{UV}}$ ($F356W$ for $z_{\text{spec}} \le 5$; $F444W$ for $z_{\text{spec}} > 5$).
Below, the Euler diagram and bar plot show that having a point source and a V-shape implies $\sim$80 percent odds of a broad line. 
Together, the RUBIES survey demonstrate that combining color-slope, compact-morphology, and spectroscopic-shape cuts can define LRD as a population \citep{hviding2025rubies}.
}
\label{fig:hviding}
\end{figure}

\begin{figure}[ht]
\begin{center}
\includegraphics[width=\textwidth]{Furtak.png}% This is a *.eps file
\end{center}
\caption{
Figure credit: Lukas Furtak and Adi Zitrin, as adapted from \citet{furtak2025investigating}, their Figure 1. 
This figure provides direct spectroscopic evidence for accretion in the little red dot A2744-QSO1. JWST/NIRSpec-prism observations of the multiply imaged source A2744-QSO1 at multiple epochs capture the full Balmer region and reveal broad H$\alpha$ and H$\beta$ and their variability. After placing the spectra on a common reference and scaling line profiles by the local continua, the broad shapes persist while only modest equivalent-width changes remain, which is consistent with AGN variability. These measurements show that some LRD host broad-line regions even when continuum variability is weak, supporting an interpretation in which at least part of the LRD population traces active black hole growth \citep{furtak2025investigating}.}
\label{fig:furtak}
\end{figure}



%%% If you don't add the figures in the LaTeX files, please upload them when submitting the article.
%%% Frontiers will add the figures at the end of the provisional pdf automatically
%%% The use of LaTeX coding to draw Diagrams/Figures/Structures should be avoided. They should be external callouts including graphics.

\end{document}
